\documentclass[10pt]{article}
\usepackage[a4paper,margin=25mm]{geometry}
\usepackage{amsmath,amssymb,amsthm}
\usepackage[hidelinks]{hyperref}
\newtheorem{proposition}{Proposition}
\newcommand{\code}[1]{\texttt{\small\detokenize{#1}}}
\title{A counterexample to Conjecture 00000006542\\
\large Ihara poles are not reciprocals of adjacency eigenvalues}
\author{Chengzhe Gao}
\date{October 3, 2026}
\begin{document}
\maketitle

\paragraph{The claim being refuted.}
The English source asserts that ``the reciprocals of the poles are the
adjacency spectrum.'' The Chinese source states the same correspondence.
In particular, it implies that if $a\ne0$ is an Ihara zeta pole, then
$a^{-1}$ is an adjacency eigenvalue. We refute this necessary implication
for the finite, connected, simple, $2$-regular graph $C_3$. The source
does not exclude cycles or impose minimum degree at least $3$.
No interpretation of its other clauses about an index, a domain, or
coverings is needed for this counterexample.

\paragraph{The standard objects.}
We use the usual Ihara definition in which a prime is a primitive,
closed, backtrackless, tailless directed walk, modulo cyclic rotation.
Reversing the direction does not identify two primes. Thus
\[
 Z_G(u)=\prod_{[P]}(1-u^{\ell(P)})^{-1}.
\]
These conventions are the ones in Horton--Stark--Terras,
Section~2, equations (2.1)--(2.2)\cite{hst}. Their definition applies to
finite connected graphs of fundamental-group rank at least one with
no degree-one vertices; $C_3$ satisfies these conditions. For $C_3$ the
product below is finite, so no convergence theorem or determinant
formula is required.

\begin{proposition}
The Ihara zeta function of $C_3$ has a pole at $1$, but $1$ is not an
eigenvalue of its adjacency matrix over $\mathbb C$.
\end{proposition}
\begin{proof}
Use vertices $0,1,2$ and edges $01,12,20$. At any directed edge, there
is exactly one next edge that is incident to its head and does not
immediately reverse it. The six directed edges therefore form the two
successor orbits
\[
 (01,12,20),\qquad (02,21,10).
\]
Every backtrackless walk stays in one of these orbits. It closes
without a tail if and only if its positive length is divisible by $3$.
A closed walk of length $3q$, with $q\ge2$, is the $q$th power of its
first three edges. Conversely a length-three closed walk cannot be a
proper power of a shorter reduced closed walk, since the latter would
also have positive length divisible by $3$. Hence all prime walks have
length $3$. Cyclic rotation identifies precisely the three starting
edges in each displayed orbit, giving exactly two prime classes.
Consequently the defining Euler product itself gives
\[
 Z_{C_3}(u)=\frac{1}{(1-u^3)^2}
          =\frac{1}{(u-1)^2(1+u+u^2)^2}.
\]
The last factor has value $9$ at $u=1$, and the numerator has value
$1$. Thus there is no cancellation, and $u=1$ is a pole of exact order
$2$. Its reciprocal is $1$.

The adjacency action on an arbitrary vector $(x,y,z)\in\mathbb C^3$ is
\[
 A(x,y,z)=(y+z,x+z,x+y),\qquad
 A=\begin{pmatrix}0&1&1\\1&0&1\\1&1&0\end{pmatrix}.
\]
If $Av=v$, then $x=y+z$, $y=x+z$, and $z=x+y$.
The first two equations imply $x=x+2z$, so $z=0$ and $x=y$.
The third now gives $2x=0$, and therefore $x=y=z=0$.
There is no nonzero eigenvector with eigenvalue $1$.
\end{proof}

The displayed failure concerns membership in the spectrum, so it does
not depend on whether the source meant a set or a multiset of poles
and eigenvalues. In the usual theory the relevant direct spectral
object is the nonbacktracking directed-edge operator; replacing it
by the adjacency operator is the error tested here.

\paragraph{Formalization and coverage.}
The accompanying \code{lean/Main.lean} uses only Lean 4.19.0 and
\code{Std}. Its finite computations concern exactly the graph and
polynomials above. The classification of walks covers arbitrary
positive lengths, not a bounded search.

The graph has vertex type \code{Fin 3}, with two vertices adjacent
exactly when distinct. The six darts have explicit tails, heads, and
reverses. Theorems \code{triangle_connected_regular} and
\code{all_oriented_edges} verify regularity, connectivity, completeness
and uniqueness of the dart encoding, and reversal. The predicate
\code{Step} is equality of incident endpoints together with absence of
immediate reversal, and \code{step_is_forced} proves its equivalence to
the stated successor permutation.

A \code{ClosedWalk} is a positive-length periodic dart sequence with
\code{Step} at every pair of consecutive positions. This condition
includes the junction between the last and first edge, and thus the
usual tailless condition. Every ordinary finite cyclic dart word of
every positive length is formally mapped into this representation by
\code{fromFiniteWord}; its periodic extension is indexed modulo that
length. Conversely, taking the first declared number of darts of a
\code{ClosedWalk} recovers such a finite cyclic word, because its
periodicity identifies the boundary positions. Thus no finite closed
walks are omitted. Primitivity explicitly excludes a shorter positive
period whose integer repetition number is at least two. Such a period
is itself a closed reduced walk, so this is the usual proper-power
condition. The symbolic theorem
\code{primitive_iff_length_three} proves the classification for all
lengths.

\code{PrimeCycle} is an actual quotient of primitive walks by rotation,
not a two-element type declared to be the primes. The theorem
\code{rotations_iff_direction} proves the equivalence relation and its
classification. The maps \code{cycleCode} and \code{codeCycle} are
proved inverse. In particular, \code{prime_cycles_complete} and
\code{prime_cycles_nodup} prove that the two displayed representatives
are the full index set of the Ihara product, without duplicates.
\code{iharaZeta} multiplies the Euler factors over this complete list.
\code{actual_euler_product} then verifies its numerator and denominator.

Polynomials are finite lists of integer coefficients in ascending
degree. Addition and multiplication use coefficient addition and full
convolution; evaluation is Horner evaluation. A rational expression
is its polynomial numerator and denominator. The predicate
\code{HasPoleAt} uses the exact algebraic criterion
\[
 D(u)=(u-a)^m Q(u),\quad m>0,\quad Q(a)\ne0,\quad N(a)\ne0.
\]
For integer-coefficient rational functions considered over $\mathbb C$,
this criterion asserts an ordinary meromorphic pole of exact order
$m$: the quotient $N/Q$ is analytic and nonzero near $a$.
The Lean theorem \code{pole_at_one_order_two} verifies this factorization
for the actual Euler product, and \code{residual_value} verifies
$Q(1)=9$. Its denominator is nonzero since it has value $1$ at $0$.
This finite algebraic pole certificate avoids assuming an unproved
analytic continuation or a graph-zeta determinant identity.

The eigenvalue predicate uses all functions \code{Fin 3} to an
arbitrary scalar type, a nonzero coordinate, and the actual matrix
action, formed by summing its three entries in each row. The statement
\code{one_not_adjacency_eigenvalue} is proved for every scalar structure
\code{ScalarLaws} with additive associativity, zero, cancellation,
no additive $2$-torsion, and the usual zero and unit multiplication
laws. These are all standard laws of $\mathbb C$ (also of any field of
characteristic different from $2$). Hence the proof applies to
arbitrary complex eigenvectors; it is not restricted to integer or
rational vectors. The separately supplied integer instance only
demonstrates consistency of the small algebraic interface.

Finally \code{conjecture6542_false} rejects the necessary implication
from this actual pole to this actual eigenvalue. The auxiliary script
checks graph walks through length $15$, the same exact polynomial
factorization, and $\det(I-A)=-4$; it is not used to justify the
all-length classification or any Lean theorem. The printed axioms of
the principal theorems are among \code{propext},
\code{Classical.choice}, and \code{Quot.sound}. There are no placeholders,
extra axioms, or unchecked native decision procedures.

\begin{thebibliography}{9}
\bibitem{hst}
M.~D.~Horton, H.~M.~Stark, and A.~A.~Terras,
\emph{What are zeta functions of graphs and what are they good for?},
Section~2, author-hosted manuscript,
\url{https://mathweb.ucsd.edu/~aterras/snowbird.pdf}.
\end{thebibliography}
\end{document}
